Once every electrode in the four curated regions had been evaluated individually, the single-electrode results were aggregated to the same five scalp regions used in the regional channel-subset analysis (Frontal, Central, Parietal, Occipital and Posterior, the last being the union of Parietal and Occipital). Each electrode's own single-electrode macro $F_1$ was averaged over sessions, and those per-electrode means were then averaged across the electrodes the region contains, exactly as in the region-mean aggregation used for the full montage and the regional subsets (Table~\ref{tab:exp1_single_channel_region}). On Internal, the resulting region means were Frontal $64.95\%$ ($n=6$), Central $12.82\%$ ($n=6$), Parietal $1.71\%$ ($n=4$), Occipital $2.00\%$ ($n=8$) and Posterior $1.90\%$ ($n=12$); on Cao2018, Frontal $62.89\%$ ($n=4$), Central $32.48\%$ ($n=6$), Parietal $15.51\%$ ($n=2$), Occipital $6.60\%$ ($n=6$) and Posterior $8.83\%$ ($n=8$). In both corpora the same anterior-to-posterior hierarchy seen at the electrode level and in the regional-subset analysis was preserved when single-electrode results were aggregated by region.

Each region's single-electrode mean was compared, session-paired, against the same electrodes' region mean scored inside the full-montage run, using the same two-tailed Wilcoxon signed-rank framework as above, Bonferroni-corrected over the 5 regions tested within each dataset (a family separate from the electrode-level family above). Every region in both corpora scored significantly lower operating as independent single-electrode detectors than as part of the full montage (all adjusted $p<0.001$; Table~\ref{tab:exp1_single_channel_region}). On Internal, the paired differences were Frontal $\Delta F_1=-4.74$ ($r_{\mathrm{rb}}=-0.80$), Central $-1.98$ ($r_{\mathrm{rb}}=-0.77$), Parietal $-0.23$ ($r_{\mathrm{rb}}=-0.75$), Occipital $-0.30$ ($r_{\mathrm{rb}}=-0.92$) and Posterior $-0.28$ ($r_{\mathrm{rb}}=-0.90$); on Cao2018, Frontal $-6.02$ ($r_{\mathrm{rb}}=-0.72$), Central $-5.65$ ($r_{\mathrm{rb}}=-0.79$), Parietal $-3.19$ ($r_{\mathrm{rb}}=-0.56$), Occipital $-1.64$ ($r_{\mathrm{rb}}=-0.69$) and Posterior $-2.03$ ($r_{\mathrm{rb}}=-0.65$). The Internal frontal region showed the largest absolute reduction of any Internal region, mirroring the electrode-level pattern in which AF3/AF4 (but not Fp1/Fp2) changed significantly; on Cao2018 the frontal reduction was likewise the largest in absolute terms, driven by the large F3/F4 effects noted above rather than by the frontopolar pair.
